\section{Pedagogical Activity}\label{pedagogical-activity}

My pedagogical activity has involved the coordination and teaching of
courses at IST and the supervision of doctoral and master students, in
the broad disciplinary areas of Electrical and Computer Engineering,
Aerospace Engineering, and Computer Science.

\subsection{Courses}\label{courses}

I have lectured more than \textbf{21 different non-tutorial curricular
units} in the 1\textsuperscript{st}, 2\textsuperscript{nd} and
3\textsuperscript{rd} cycles, most on the core of the Integrated Master
Programme in Electrical and Computer Engineering, on the disciplinary
areas of Control (Control Theory, Modeling and Simulation, Systems
Theory, among others), Robotics (Automation, Autonomous Systems,
Robotics, among others), and Signals (Signal Processing, Circuit Theory,
Machine Learning, Algorithms, Computer Vision, Neural Networks, among
others).

In the last years I have been responsible for the MSc courses of
Distributed Real Time Control Systems (\textasciitilde80 students) and
Machine Learning (\textasciitilde270 students). In the last two years I
have become responsible for the PhD course of Estimation and
Classification (\textasciitilde10 students). I was granted three times
the IST Excellent Teaching award, issued by the IST Course Unit Quality
System (based on students' surveys), on the course of Distributed Real
Time Control Systems (2016/2017, 2020/2021 and 2023/2024). The full list
of courses is provided in Appendix \ref{courses-1}, including the ranks given by
students to my teaching performance, according to the student's surveys
organized by IST pedagogical council program QUC (IST Course Unit
Quality System).

\subsection{Innovation and Pedagogical
Materials}\label{innovation-and-pedagogical-materials}

My teaching activities are characterized by continued efforts to
introduce new contents in the courses I have lectured, both in the
theoretical and practical components, with emphasis in experimental and
laboratory resources. Full lists of my innovation actions and
pedagogical materials are presented in Appendices \ref{innovation} and \ref{pedagogical-materials}.

I started my teaching career in 1995 as a Teaching Assistant associated
to the laboratorial component of courses of scientific area of Systems
Decision and Control, where I always thrived for the introduction of new
concepts, projects and equipment. In 1996 I prepared, together with the
\textbf{Neural Networks} course coordinator, Prof. Luís Borges de
Almeida, a set of original lab projects involving new software for the
design and training of neural networks. In 1998 I prepared, together
with the \textbf{Computer Vision} course coordinator, Prof. José
Santos-Victor, an original lab project for 3D reconstruction based on
cast shadows. In 2002 I prepared, together with the \textbf{Modeling,
Identification and Control} course coordinator, Prof. João Lemos, a new
project for the control of a ball-in-beam educational kit. In the follow
up of this project, during 2003-2007 I collaborated in several proposals
for the improvement of the control labs with Quanser educational kits
and I/O acquisition boards to be used in the teaching of control topics
in several courses, corresponding to an investment around 50kEur.

In 2003, I got the coordination of the course on \textbf{Modeling,
Identification and Control} (2003-2011), where I developed exciting lab
projects for the Quanser educational kits ball-in-beam and flexible-bar.
Simultaneously, I prepared, together with the Optimal and Adaptive
Control course coordinator, Prof. João Lemos, a lab project involving
the implementation of the adaptive controller MUSMAR in Simulink
Real-Time, for the control of a DC motor. In 2003 I prepared, together
with the \textbf{Automation and Industrial Processes} course
coordinator, Prof. Paulo Oliveira, a lab project involving a new
automation cell composed of conveyor belts and robot manipulators.

Since 2007, I have been coordinating the 2\textsuperscript{nd} cycle
course on \textbf{Distributed Real Time Control Systems} (with
interruption in 2010/2011, 2011/2012, and 2014/2015). This is a
projects-based course where students develop a practical control system
involving knowledge on control theory, system modeling and
identification, distributed optimization and control, real-time
concurrent systems, and control communication networks. I extensively
reformulated the course by introducing decentralized control and
distributed optimization topics. Also, I designed a new project that
students can implement with accessible materials (basic electronic,
communication and computational elements) and take home for allowing
autonomous work. In 2013 I prepared a take-home kit that allows students
to develop their project anywhere/anytime with real hardware. Since
then, students can manage their own time and dedication to the project,
which promotes their learning skills. During the years of the COVID19
pandemics, the project developed without disruptions, as it was already
fully prepared for in-home execution.

During the year of 2009/2010 I coordinated the \textbf{Machine Learning}
course, during the sabbatical leave of Prof. Luis Borges de Almeida. As
I was involved from the beginning of the course in the preparation of
laboratory projects, and taught the laboratory classes of the previous
year, I was invited by the usual course coordinator to replace him on
his leave.

The idea of take-home projects initiated by me on \textbf{Distributed
Real Time Control Systems} course, and the disruptions in the teaching
of project-based and lab-based courses that emerged from the pandemic
crisis, were on the genesis of three Pedagogical Innovation Projects in
the years 2021 and 2022, that allowed not only the students of
\textbf{Distributed Real Time Control Systems} but also of
\textbf{Control Theory}, and \textbf{Autonomous Systems}, to take the
hardware and instrumentation of their projects home and improve their
autonomy, creativity, and time management skills.

\subsection{Supervision}\label{supervision}

I have supervised and co-supervised \textbf{22 concluded PhD students},
\textbf{10 ongoing PhD students}, and more than \textbf{210 MSc
students} (concluded).

Many of my PhD students are joint supervisions with international
partners in dual-degree programmes: 3 in the IST-CMU programme, 3 in the
IST-EPFL programme, and 1 in co-tutelle with UPMC, Paris, France. My
supervised PhD student Pedro Vicente was awarded \textbf{Best National
PhD Thesis in Robotics Award of 2022}, Portuguese Society of Robotics,
and the \textbf{2\textsuperscript{nd} Place of the Best Iberian PhD
Thesis in Robotics 2022}, Spanish Society for Research and Development
in Robotics and the Portuguese Robotics Society. My supervised PhD
student Rui Figueiredo was awarded the \textbf{Honorable Mention for the
Best National PhD Thesis Award of 2019/2020}, Portuguese Association of
Pattern Recognition. My co-supervised PhD student Ricardo Cabral was
awarded the \textbf{IBM Prize 2014}.

Most of my MSc students are from the courses on the responsibility of
the Department of Electrical and Computer Engineering at Técnico (MEEC,
MEE, MEAer) but I also have supervised students from courses in other
Departments (Physics, Mathematics, Computer Science, Biomedical
Engineering, Civil Engineering), other national universities (FCT/UNL)
and other international universities (HTW Berlin, Germany), (TU/e,
Netherlands). One of my MSc students (Filipe Veiga) was awarded the
\textbf{Best MSc Thesis Award} by the Portuguese Society of Robotics
(SPR) in 2012. Another of my MSc students (João Castelo-Branco) was
awarded one of the \textbf{Outstanding Master Thesis diplomas} of 2026
by DEEC.

Beyond supervision of students in formal educational programmes, I
supervised the work of \textbf{9 post-doctoral researchers}, \textbf{84
research assistants}, \textbf{23 visiting researchers}.

The full list of my MSc students, PhD students, post-doctoral
researchers, research assistants and visiting researchers is provided in
Appendix \ref{supervision-activities}.
